

\documentclass{article}

\usepackage{fancyhdr}
\usepackage{extramarks}
\usepackage{amsmath}
\usepackage{amsthm}
\usepackage{amsfonts}
\usepackage{tikz}
\usepackage[plain]{algorithm}
\usepackage{algpseudocode}
\usepackage{amsmath,amssymb}

\usetikzlibrary{automata,positioning}

%
% Basic Document Settings
%

\topmargin=-0.45in
\evensidemargin=0in
\oddsidemargin=0in
\textwidth=6.5in
\textheight=9.0in
\headsep=0.25in

\linespread{1.1}

\pagestyle{fancy}
\lhead{\hmwkAuthorName}
\chead{\hmwkClass: \hmwkTitle}
\rhead{}
\lfoot{\lastxmark}
\cfoot{\thepage}

\renewcommand\headrulewidth{0.4pt}
\renewcommand\footrulewidth{0.4pt}

\setlength\parindent{0pt}

%
% Create Problem Sections
%

\newcommand{\enterProblemHeader}[1]{
    \nobreak\extramarks{}{Problem \arabic{#1} continued on next page\ldots}\nobreak{}
    \nobreak\extramarks{Problem \arabic{#1} (continued)}{Problem \arabic{#1} continued on next page\ldots}\nobreak{}
}

\newcommand{\exitProblemHeader}[1]{
    \nobreak\extramarks{Problem \arabic{#1} (continued)}{Problem \arabic{#1} continued on next page\ldots}\nobreak{}
    \stepcounter{#1}
    \nobreak\extramarks{Problem \arabic{#1}}{}\nobreak{}
}

\setcounter{secnumdepth}{0}
\newcounter{partCounter}
\newcounter{homeworkProblemCounter}
\setcounter{homeworkProblemCounter}{1}
\nobreak\extramarks{Problem \arabic{homeworkProblemCounter}}{}\nobreak{}



% Define theorem environment
\newtheorem*{theorem}{Theorem}

%
% Homework Problem Environment
%
% This environment takes an optional argument. When given, it will adjust the
% problem counter. This is useful for when the problems given for your
% assignment aren't sequential. See the last 3 problems of this template for an
% example.
%
\newenvironment{homeworkProblem}[1][-1]{
    \ifnum#1>0
        \setcounter{homeworkProblemCounter}{#1}
    \fi
    \section{Problem \arabic{homeworkProblemCounter}}
    \setcounter{partCounter}{1}
    \enterProblemHeader{homeworkProblemCounter}
}{
    \exitProblemHeader{homeworkProblemCounter}
}

%
% Homework Details
%   - Title
%   - Due date
%   - Class
%   - Section/Time
%   - Instructor
%   - Author
%


\newcommand{\hmwkTitle}{Problem set\ 11}
\newcommand{\hmwkDueDate}{November 24, 2025}
\newcommand{\hmwkClass}{Number Theory}
\newcommand{\hmwkClassTime}{Section 2}
\newcommand{\hmwkClassInstructor}{Dr. Eleanor McSpirit}
\newcommand{\hmwkAuthorName}{\textbf{Joseph Quinn}}

%
% Title Page
%

\title{
    \vspace{2in}
    \textmd{\textbf{\hmwkClass:\ \hmwkTitle}}\\
    \normalsize\vspace{0.1in}\small{\hmwkDueDate}\\
    \vspace{0.1in}\large{\textit{\hmwkClassInstructor\ \hmwkClassTime}}
    \vspace{3in}
}

\author{\hmwkAuthorName}
\date{}

\renewcommand{\part}[1]{\textbf{\large Part \Alph{partCounter}}\stepcounter{partCounter}\\}

%
% Various Helper Commands
%

% Useful for algorithms
\newcommand{\alg}[1]{\textsc{\bfseries \footnotesize #1}}

% For derivatives
\newcommand{\deriv}[1]{\frac{\mathrm{d}}{\mathrm{d}x} (#1)}

% For partial derivatives
\newcommand{\pderiv}[2]{\frac{\partial}{\partial #1} (#2)}

% Integral dx
\newcommand{\dx}{\mathrm{d}x}

% Alias for the Solution section header
\newcommand{\solution}{\textbf{\large Solution}}

% Probability commands: Expectation, Variance, Covariance, Bias
\newcommand{\E}{\mathrm{E}}
\newcommand{\Var}{\mathrm{Var}}
\newcommand{\Cov}{\mathrm{Cov}}
\newcommand{\Bias}{\mathrm{Bias}}

%  proof-step macro:
\newcommand{\step}[2]{& #1 & & \text{#2} \\}

\begin{document}

\maketitle
\pagebreak

\begin{homeworkProblem}
	\begin{proof}
		We prove the theorem by proving its contrapositive.
		Assume that $p$ is a prime such that
		\[
			p \ne 2 \quad\text{and}\quad p \not\equiv 1 \pmod{4}.
		\]
		Since every odd prime is congruent to either $1$ or $3$ modulo $4$, this means that
		\[
			p \equiv 3 \pmod{4}.
		\]

		We show that such a prime cannot be written as a sum of two squares.
		Suppose, for contradiction, that there exist natural numbers $a$ and $b$ such that
		\[
			p = a^{2} + b^{2}.
		\]

		Consider this equation modulo $4$.
		Since every square is congruent to either $0$ or $1$ modulo $4$, the possible values of
		$a^{2} + b^{2} \pmod{4}$ are
		\[
			0+0,\; 0+1,\; 1+0,\; 1+1,
		\]
		which yield the residues $0,1,2 \pmod{4}$.
		In particular, $a^{2} + b^{2}$ can never be congruent to $3 \pmod{4}$.

		But $p \equiv 3 \pmod{4}$, so this is impossible.
		Thus no such $a$ and $b$ can exist, and a prime congruent to $3 \pmod{4}$ cannot be a sum of two natural squares.
	\end{proof}
\end{homeworkProblem}

\begin{homeworkProblem}
	\begin{proof}
		Since $(a,b,c)$ is a primitive Pythagorean triple with $a$ even, we have
		\[
			a^2 + b^2 = c^2, \qquad \gcd(a,b,c)=1,
		\]
		and $b$ and $c$ are odd.

		\medskip
		\noindent\textbf{(i)}
		A square modulo $3$ is congruent only to $0$ or $1$. Thus
		\[
			a^2 + b^2 \equiv 0,1,2 \pmod{3}.
		\]
		The value $2 \pmod{3}$ is impossible for a square, so we must avoid the case
		$a^2 \equiv b^2 \equiv 1 \pmod{3}$, which would give $c^2 \equiv 2 \pmod{3}$.
		Therefore, not both of $a$ and $b$ are congruent to $1 \pmod{3}$.

		If both were divisible by $3$, then $3 \mid a$ and $3 \mid b$, and hence
		$3 \mid c$, contradicting primitivity.
		Thus exactly one of $a$ and $b$ is divisible by $3$.

		\medskip
		\noindent\textbf{(ii)}
		Since $a$ is even and $(a,b,c)$ is primitive, $b$ and $c$ are odd.
		Squares modulo $4$ satisfy
		\[
			x^2 \equiv
			\begin{cases}
				0 \pmod{4}, & x \text{ even}, \\
				1 \pmod{4}, & x \text{ odd}.
			\end{cases}
		\]
		Thus
		\[
			a^2 \equiv 0 \pmod{4}, \qquad b^2 \equiv 1 \pmod{4}, \qquad
			c^2 = a^2 + b^2 \equiv 1 \pmod{4}.
		\]

		Since $b$ is odd, $4 \nmid b$.
		If $4 \mid a$ and $4 \mid b$, then both would be even, contradicting primitivity.
		Thus exactly one of $a$ and $b$ is divisible by $4$ (in fact $a$ is).

		\pagebreak
		\textbf{(iii)}
		Squares modulo $5$ take the values $0,1,4$.
		Hence $a^2 + b^2 \pmod{5}$ may take the values
		\[
			0,1,4,2,3.
		\]
		But $c^2$ must be $0,1$, or $4 \pmod{5}$.
		Thus the cases $a^2 + b^2 \equiv 2,3 \pmod{5}$ must be avoided.

		If neither $a$ nor $b$ is divisible by $5$, then $a^2,b^2\in\{1,4\}$.
		To avoid the forbidden residues $2$ and $3$, we must have one of $a^2,b^2$
		congruent to $1$ and the other to $4$, giving
		\[
			c^2 \equiv 0 \pmod{5},
		\]
		so $5 \mid c$.

		If exactly one of $a,b$ is divisible by $5$, say $5 \mid a$, then
		\[
			a^2 \equiv 0,\qquad b^2 \equiv 1 \text{ or } 4,
		\]
		so $c^2 \equiv 1$ or $4 \pmod{5}$, and hence $5 \nmid c$.

		If both $a$ and $b$ were divisible by $5$, primitivity would be violated.

		Therefore exactly one of $a$, $b$, and $c$ is divisible by $5$.
	\end{proof}

\end{homeworkProblem}

\begin{homeworkProblem}

	To find infinitely many integer solutions to
	\[
		x^{2} + 2y^{2} = z^{2},
	\]
	we imitate the construction of Pythagorean triples.
	Let $m$ and $n$ be natural numbers with $m > n$.
	Consider the triple
	\[
		(x,y,z) = (m^{2} - 2n^{2},\; 2mn,\; m^{2} + 2n^{2}).
	\]

	We verify that this satisfies the equation.
	Compute the left-hand side:
	\[
		(m^{2} - 2n^{2})^{2} + 2(2mn)^{2}
		= m^{4} - 4m^{2}n^{2} + 4n^{4} + 8m^{2}n^{2}
		= m^{4} + 4m^{2}n^{2} + 4n^{4}.
	\]
	The right-hand side is
	\[
		(m^{2} + 2n^{2})^{2}
		= m^{4} + 4m^{2}n^{2} + 4n^{4}.
	\]

	Thus
	\[
		(m^{2} - 2n^{2})^{2} + 2(2mn)^{2}
		= (m^{2} + 2n^{2})^{2},
	\]
	so $(x,y,z)$ is indeed a solution.

	Since every choice of $m>n$ produces a different triple, this gives infinitely many integer solutions to
	\[
		x^{2} + 2y^{2} = z^{2}.
	\]

\end{homeworkProblem}

\begin{homeworkProblem}

	We determine whether the Diophantine equation
	\[
		x^{2} + y^{2} = 175
	\]
	has an integer solution.

	Factor the right-hand side:
	\[
		175 = 25 \cdot 7 = 5^{2} \cdot 7.
	\]
	The prime $7$ satisfies $7 \equiv 3 \pmod{4}$ and appears with exponent $1$ in the
	factorization of $175$.

	By the sum of two squares theorem, a natural number can be written as a sum of two
	squares if and only if every prime $q \equiv 3 \pmod{4}$ in its prime factorization
	occurs with an even exponent. Since $7$ occurs to an odd power, $175$ cannot be written
	as a sum of two squares.

	Therefore, the equation $x^{2} + y^{2} = 175$ has no integer solutions.

\end{homeworkProblem}

\begin{homeworkProblem}

	Let $x,y,z \in \mathbb{Z}$ and let $p$ be a prime.

	\medskip
	\noindent\textbf{(i)}
	Assume that
	\[
		x^{p-1} + y^{p-1} = z^{p-1}.
	\]
	Suppose, for contradiction, that $p \nmid x$, $p \nmid y$, and $p \nmid z$.
	Then by Fermat's Little Theorem,
	\[
		x^{p-1} \equiv y^{p-1} \equiv z^{p-1} \equiv 1 \pmod{p}.
	\]
	Reducing the given equation modulo $p$ yields
	\[
		1 + 1 \equiv 1 \pmod{p},
	\]
	which simplifies to $2 \equiv 1 \pmod{p}$, an impossibility.
	Hence at least one of $x,y,z$ must be divisible by $p$.
	Therefore
	\[
		p \mid xyz.
	\]

	\medskip
	\noindent\textbf{(ii)}
	Assume that
	\[
		x^{p} + y^{p} = z^{p}.
	\]
	Applying Fermat's Little Theorem gives
	\[
		x^{p} \equiv x,\qquad y^{p} \equiv y,\qquad z^{p} \equiv z \pmod{p}.
	\]
	Reducing the equation modulo $p$ yields
	\[
		x + y \equiv z \pmod{p},
	\]
	or equivalently,
	\[
		x + y - z \equiv 0 \pmod{p}.
	\]
	Thus
	\[
		p \mid (x + y - z).
	\]

\end{homeworkProblem}

\end{document}
